\section{上同调与基变换}\label{src:A.6}
\begin{epigraph}
几乎没有哪种理论比 [线性代数] 更初等, 尽管一代代教授和教材作者用荒唐的矩阵计算掩盖了它的简单性.
\epigraphauthor{Jean Dieudonné}
\end{epigraph}
若 $f:X\to Y$ 为 Noether 概形之间的固有态射, $F$ 为 $X$ 上的凝聚层, 则上同调的有限性 (\bookref{A.5.3}) 意味着 $\mathrm R^i f_*F$ 凝聚. 我们常常还想知道:
\begin{enumerate}[(a)]
\item $\mathrm R^i f_*F$ 何时是 $Y$ 上的向量丛?
\item $\mathrm R^i f_*F$ 的构造何时与基变换交换? 即对诱导 Cartesian 图
\[
\begin{tikzcd}[column sep=3em,row sep=2em]
X_{Y'}\arrow[r,"g'"]\arrow[d,"f'"']&X\arrow[d,"f"]\\
Y'\arrow[r,"g"']&Y,
\end{tikzcd}
\]
的概形映射 $g:Y'\to Y$, 比较映射
\begin{equation}\label{src:A.6.1}
\phi_{Y'}^i:g^*\mathrm R^if_*F\longrightarrow\mathrm R^if'_*g'^*F
\end{equation}
何时为同构?
\end{enumerate}
当 $f:X\to Y$ 平坦时, 平坦基变换 (\bookref{A.2.12}) 告诉我们 \bookeqref{A.6.1} 总是同构. 上同调与基变换 (\bookref{A.6.8}) 则在 $F$ 在 $Y$ 上平坦时给出答案.

上同调与基变换是模理论的重要工具. 可以用它验证对象族的性质, 并在模空间上构造向量丛. 例如, 对光滑曲线族 $\pi:C\to S$, 可以验证对 $k>0$, $\pi_*(\Omega_{C/S}^{\otimes k})$ 是构造与 $S$ 上基变换交换的向量丛, 并证明当 $k\geq3$ 时, $C$ 典范地嵌入 $\mathbb P(\pi_*(\Omega_{C/S}^{\otimes k}))$ (命题 \bookref{6.1.16}). 我们在研究 $\mathcal M_g$ 时应用上同调与基变换, 用它验证 $\mathcal M_g$ 为代数叠 (定理 \bookref{4.1.17}), 并建立各种几何性质. 对万有族 $\pi:\mathcal U_g\to\mathcal M_g$ 应用上同调与基变换, 可知 $\pi_*(\Omega_{\mathcal U_g/\mathcal M_g})$ 为 $\mathcal M_g$ 上秩 $g$ 的向量丛, 称为 Hodge 丛 (例 \bookref{5.1.4}).

\subsection{上同调与基变换的各种表述}\label{src:A.6.1-subsection}
先给出上同调与基变换的关键代数版本, 其他版本都由它建立.
\begin{theorem}[上同调与基变换 I]\label{src:A.6.2}
设 $X\to\Spec A$ 为 Noether 概形之间的固有态射, $F$ 为 $X$ 上在 $A$ 上平坦的凝聚层. 存在有限且局部自由 $A$-模的复形
\[
K^\bullet:0\longrightarrow K^0\longrightarrow K^1\longrightarrow\cdots\longrightarrow K^n\longrightarrow0,
\]
使得对所有 $i$ 有 $\mathrm H^i(X,F)=\mathrm H^i(K^\bullet)$. 此外, 对每个 $A$-模 $M$, 有 $\mathrm H^i(X,F\otimes_A M)=\mathrm H^i(K^\bullet\otimes_A M)$. 特别地, 对概形态射 $\Spec B\to\Spec A$, 有 $\mathrm H^i(X_B,F_B)=\mathrm H^i(K^\bullet\otimes_A B)$, 其中 $X_B:=X\times_{\Spec A}\Spec B$, $F_B$ 为 $F$ 到 $X_B$ 的拉回.
\end{theorem}
\begin{proof}
选取 $X$ 的有限仿射覆盖 $\{U_i\}$, 并考虑以 $F$ 为系数的相应交错 Čech 复形 $C^\bullet$. 则 $C^\bullet$ 为平坦 (但非有限生成) $A$-模的有限复形, 且 $\mathrm H^i(X,F)=\mathrm H^i(C^\bullet)$. 归纳地改造 $C^\bullet$, 构造与其拟同构的有限平坦 $A$-模的有限复形 $K^\bullet$, 即得结论. 见 \cite[第 $46$ 页的定理]{Mum70a}, \cite[Tag 07VK]{SP} 和 \cite[$28.2.1$]{Vak17}.
\end{proof}
\begin{remark}[完美复形]\label{src:A.6.3}
Noether 概形 $X$ 上凝聚层的有界复形 $K^\bullet$ 称为\newterm{完美复形}{perfect complex}, 如果存在仿射覆盖 $X=\bigcup_i U_i$, 使每个 $K^\bullet|_{U_i}$ 都拟同构于 $U_i$ 上向量丛的有界复形. 若 $X$ 仿射 (或更一般地具有消解性质, 即每个凝聚层都是向量丛的商), 则 $K^\bullet$ 完美, 当且仅当它拟同构于 $X$ 上向量丛的有界复形 \cite[Tag 0F8F]{SP}. 此外, $D_{\mathrm{QCoh}}(X)$ 中的紧对象恰好是完美复形 \cite[Tag 09M8]{SP}.

采用这一术语, 定理 \bookref{A.6.2} 可以表述为: $\mathrm Rf_*F\in D_{\mathrm{Coh}}^b(\Spec A)$ 完美. 更一般地, 若 $F^\bullet$ 为 $X$ 上的完美复形, 则 $\mathrm Rf_*F^\bullet$ 也完美 \cite[Tag 0A1H]{SP}.
\end{remark}
定理 \bookref{A.6.2} 告诉我们, 上同调 $\mathrm H^i(X,F)$ 可以作为 $\Spec A$ 上向量丛的有界复形 $K^\bullet$ 的上同调来计算, 从而把上同调问题归结为线性代数.
\begin{theorem}[半连续性定理]\label{src:A.6.4}
设 $X\to Y$ 为 Noether 概形之间的固有态射, $F$ 为 $X$ 上在 $Y$ 上平坦的凝聚层.
\begin{enumerate}[(1)]
\item 对每个 $i\geq0$, 函数
\[
Y\longrightarrow\mathbb Z,\qquad y\longmapsto\mathrm h^i(X_y,F_y)
\]
上半连续.
\item 函数
\[
Y\longrightarrow\mathbb Z,\qquad y\longmapsto\chi(X_y,F_y)=\sum_{i=0}^\infty(-1)^i\mathrm h^i(X_y,F_y)
\]
局部常值.
\end{enumerate}
\end{theorem}
\begin{proof}
可假设 $Y=\Spec A$, 从而应用定理 \bookref{A.6.2}: 存在有限且局部自由 $A$-模的有界复形 $K^\bullet:\cdots\to K^{i-1}\xrightarrow{d^{i-1}}K^i\xrightarrow{d^i}K^{i+1}\to\cdots$, 使得对所有 $y\in Y$, 有 $\mathrm H^i(X_y,F_y)=\mathrm H^i(K^\bullet\otimes_A\kappa(y))$. 使用 $\im(d^i\otimes\kappa(y))=(K^i\otimes_A\kappa(y))/\ker(d^i\otimes\kappa(y))$, 得到
\begin{equation}\label{src:A.6.5}
\begin{aligned}
\mathrm h^i(X_y,F_y)&=\dim_{\kappa(y)}\ker(d^i\otimes\kappa(y))-\dim_{\kappa(y)}\im(d^{i-1}\otimes\kappa(y))\\
&=\dim_{\kappa(y)}K^i\otimes\kappa(y)-\dim_{\kappa(y)}\im(d^i\otimes\kappa(y))\\
&\qquad-\dim_{\kappa(y)}\im(d^{i-1}\otimes\kappa(y)).
\end{aligned}
\end{equation}
由于 $\dim_{\kappa(y)}\im(d^i\otimes\kappa(y))$ 和 $\dim_{\kappa(y)}\im(d^i\otimes\kappa(y))$ 都下半连续, 结论成立. 另见 \cite[第 $47$ 页]{Mum70a}, \cite[定理 $12.8$]{Har77} 或 \cite[定理 $25.1.1$]{Vak17}.
\end{proof}
为证明更强的结果, 需要更精细的线性代数. 我们遵循 Eric Larson 的一个论证, 它也见于 \cite[\S$25.2$]{Vak17}.
\begin{proposition}\label{src:A.6.6}
设 $X$ 为概形, $\phi:E\to F$ 为 $X$ 上秩分别为 $e$ 和 $f$ 的向量丛的映射. 对每个点 $x\in X$ 和整数 $r\leq\min(e,f)$, 下列条件等价:
\begin{enumerate}[(1)]
\item $\operatorname{coker}(\phi)$ 在 $x$ 的某个开邻域中为秩 $f-r$ 的向量丛;
\item 存在 $x$ 的开邻域 $U$ 和等同 $E|_U\simeq\mathcal O_U^{\oplus e}$, $F|_U\simeq\mathcal O_U^{\oplus f}$, 使得 $\phi|_U$ 对应于投影 $\mathcal O_U^{\oplus e}\to\mathcal O_U^{\oplus r}$ 与到前 $r$ 个直和项的包含 $\mathcal O_U^{\oplus r}\hookrightarrow\mathcal O_U^{\oplus f}$ 的复合;
\item $\ker(\phi)\otimes\kappa(x)\to\ker(\phi\otimes\kappa(x))$ 满射.
\end{enumerate}
若 $X$ 还约化, 则上述条件等价于
\begin{enumerate}[(1)]\setcounter{enumi}{3}
\item 存在 $x$ 的开邻域 $U$, 使得对所有 $u\in U$, $\phi\otimes\kappa(u)$ 的秩为 $r$.
\end{enumerate}
若这些条件成立, 则 $\operatorname{coker}(\phi)$, $\ker(\phi)$ 和 $\im(\phi)$ 都是向量丛, 且它们的构造与基变换交换.
\end{proposition}
\begin{proof}
这是一个线性代数练习.
\end{proof}
若上述等价条件成立, 则称 $\phi:E\to F$ 为\textit{向量丛的映射} (秩为 $r$). 现在可以给出 Grauert 定理的简短证明.
\begin{theorem}[Grauert 定理]\label{src:A.6.7}
设 $f:X\to Y$ 为 Noether 概形之间的固有态射, $Y$ 约化且连通. 设 $F$ 为 $X$ 上在 $Y$ 上平坦的凝聚层. 对每个整数 $i$, 下列条件等价:
\begin{enumerate}[(1)]
\item 函数 $y\mapsto\mathrm h^i(X_y,F_y)$ 为常值;
\item $\mathrm R^if_*F$ 为向量丛, 且比较映射
\[
\phi_y^i:\mathrm R^if_*F\otimes\kappa(y)\longrightarrow\mathrm H^i(X_y,F_y)
\]
对所有 $y\in Y$ 都是同构.
\end{enumerate}
若这些条件成立, 则 $\mathrm R^if_*F$ 的构造与任意映射 $T\to Y$ 的基变换交换.
\end{theorem}
\begin{proof}
方向 (1) $\Rightarrow$ (2) 显然. 对于逆向, 由于问题是 Zariski 局部的, 可归结为 $Y=\Spec A$. 令
\[
K^\bullet:\cdots\longrightarrow K^{i-1}\xrightarrow{d^{i-1}}K^i\xrightarrow{d^i}K^{i+1}\longrightarrow\cdots
\]
为定理 \bookref{A.6.2} 给出的 $Y$ 上向量丛的复形. 由于 $y\mapsto\mathrm h^i(X_y,F_y)$ 为常值, 恒等式 \bookeqref{A.6.5} 意味着 $y\mapsto\dim_{\kappa(y)}\im(d^{i-1}\otimes\kappa(y))$ 和 $y\mapsto\dim_{\kappa(y)}\im(d^i\otimes\kappa(y))$ 也为常值. 由于 $Y$ 约化, 命题 \bookref{A.6.6} 意味着 $d^{i-1}$ 和 $d^i$ 为向量丛的映射, $\im(d^i)$ 和 $\operatorname{coker}(d^{i-1})$ 为向量丛, 且 $\ker(d^i)$ 与基变换交换. 上同调 $\mathrm H^i(K^\bullet)=\im(d^{i-1})/\ker(d^i)$ 位于短正合列
\[
0\longrightarrow\mathrm H^i(K^\bullet)\longrightarrow\underbrace{K^i/\ker(d^i)}_{\im(d^i)}\longrightarrow\operatorname{coker}(d^{i-1})\longrightarrow0
\]
中, 因而 $\mathrm H^i(K^\bullet)$ 也是向量丛. 由于余核总与基变换交换, $\mathrm H^i(K^\bullet)=\operatorname{coker}(K^{i-1}\to\ker(d^i))$ 也如此. 另见 \cite[推论 $2$, 第 $48$ 页]{Mum70a}, \cite[推论 $12.9$]{Har77} 或 \cite[$28.1.5$]{Vak17}.
\end{proof}
Grauert 定理中约化性的假设在模理论的应用中限制很大, 因为我们常常需要在任意基底上建立对象族的性质. 幸运的是, 只需多一点线性代数, 就可以建立适用于任意基底的以下判别法.
\begin{theorem}[上同调与基变换 II]\label{src:A.6.8}
设 $f:X\to Y$ 为概形之间的固有有限表现态射, $F$ 为 $X$ 上在 $Y$ 上平坦的有限表现拟凝聚层. 假设对点 $y\in Y$ 和整数 $i$, 比较映射 $\phi_y^i:\mathrm R^if_*F\otimes\kappa(y)\to\mathrm H^i(X_y,F_y)$ 满射. 则有:
\begin{enumerate}[(a)]
\item 存在 $y$ 的开邻域 $V\subseteq Y$, 使得对每个概形态射 $Y'\to V$, 比较映射 $\phi_{Y'}^i:g^*\mathrm R^if_*F\to\mathrm R^if'_*g'^*F$ 为同构. 特别地, $\phi_y^i$ 为同构.
\item $\phi_y^{i-1}$ 满射, 当且仅当 $\mathrm R^if_*F$ 在 $y$ 的某个开邻域中为向量丛.
\end{enumerate}
\end{theorem}
\begin{remark}\label{src:A.6.9}
当 $i=0$ 时, $\phi_y^{-1}$ 是零群之间的零映射, 因而尤其是满射.
\end{remark}
\begin{proof}
先假设 $Y$ 为 Noether 概形, 并归结为仿射情形 $Y=\Spec A$. 定理 \bookref{A.6.2} 构造复形 $K^\bullet$, 使得对每个 $y\in Y$, 有复形态射
\[
\begin{tikzcd}[column sep=1.4em,row sep=2.5em]
K^{i-1}\arrow[r,"d^{i-1}"]\arrow[d]&K^i\arrow[r,"d^i"]\arrow[d]&K^{i+1}\arrow[d]&\text{计算 }\mathrm H^i(X,F)\\
K^{i-1}\otimes_A\kappa(y)\arrow[r,"d^{i-1}\otimes\kappa(y)"']&K^i\otimes_A\kappa(y)\arrow[r,"d^i\otimes\kappa(y)"']&K^{i+1}\otimes_A\kappa(y)&\text{计算 }\mathrm H^i(X_y,F_y)
\end{tikzcd}
\]
以下断言将在证明中使用两次: $\mathrm H^i(X,F)\to\mathrm H^i(X_y,F_y)$ 满射, 等价于 $\ker(d^i)\to\ker(d^i\otimes\kappa(y))$ 满射. 由于 $K^{i-1}\to K^{i-1}\otimes_A\kappa(y)$ 满射, $\im(d^{i-1})\to\im(d^{i-1}\otimes\kappa(y))$ 也满射. 对下图应用蛇引理即得断言:
\[
\begin{tikzcd}[column sep=1em,row sep=2em]
0\arrow[r]&\im(d^{i-1})\arrow[r]\arrow[d,two heads]&\ker(d^i)\arrow[r]\arrow[d]&\mathrm H^i(X,F)\arrow[r]\arrow[d]&0\\
0\arrow[r]&\im(d^{i-1}\otimes\kappa(y))\arrow[r]&\ker(d^i\otimes\kappa(y))\arrow[r]&\mathrm H^i(X_y,F_y)\arrow[r]&0
\end{tikzcd}
\]
对于 (a), 假设 $\mathrm H^i(X,F)\to\mathrm H^i(X_y,F_y)$ 满射, 由断言得 $\ker(d^i)\to\ker(d^i\otimes\kappa(y))$ 满射. 命题 \bookref{A.6.6} 意味着 $d^i:K^i\to K^{i+1}$ 为向量丛的映射; 把 $Y$ 换成 $y$ 的开邻域后, $\ker(d^i)$ 是构造与基变换交换的向量丛. 因而 $\mathrm H^i(X,F)=\operatorname{coker}(K^{i-1}\to\ker(d^i))$ 也与基变换交换. 对于 (b), 使用以下等价关系:
\begin{align*}
&\mathrm H^i(X,F)\text{ 为向量丛}\\
&\quad\overset{\bookref{A.6.6}}{\Longleftrightarrow}K^{i-1}\to\ker(d^i)\text{ 为向量丛的映射}\\
&\quad\Longleftrightarrow K^{i-1}\to K^i\text{ 为向量丛的映射}\\
&\quad\overset{\bookref{A.6.6}}{\Longleftrightarrow}\ker(d^{i-1})\otimes\kappa(y)\to\ker(d^{i-1}\otimes\kappa(y))\text{ 满射}\\
&\quad\overset{\text{断言}}{\Longleftrightarrow}\mathrm H^{i-1}(X,F)\to\mathrm H^{i-1}(X_y,F_y)\text{ 满射}.
\end{align*}
第一个等价关系由命题 \bookref{A.6.6} 得出, 因为 $\mathrm H^i(X,F)$ 为 $K^{i-1}\to\ker(d^i)$ 的余核. 第二个由以下观察得出: 由于 $d^i$ 为向量丛的映射, $\ker(d^i)$ 和 $\im(d^i)$ 为向量丛, 而映射 $\ker(d^i)\to K^i$ (其余核为 $\im(d^i)$) 也是向量丛的映射. 第三个也由命题 \bookref{A.6.6} 得出, 第四个则是证明开头的断言.

用概形的极限 (\S\bookref{B.3}) 的方法, 不难从 Noether 版本推出一般结论. 假设 $Y$ 仿射, 把 $Y=\varprojlim_{\lambda\in\Lambda}Y_\lambda$ 写成 $\mathbb Z$ 上有限型仿射概形的极限. 由于 $X\to Y$ 有限表现, 存在指标 $0\in\Lambda$ 和有限表现态射 $X_0\to Y_0$, 使得 $X\simeq X_0\times_{Y_0}Y$ (命题 \bookref{B.3.3}). 对每个 $\lambda>0$, 定义 $X_\lambda=X_0\times_{Y_0}Y_\lambda$, 并有 $X\simeq X_\lambda\times_{Y_\lambda}Y$. 由命题 \bookref{B.3.7}, 对 $\lambda\gg0$, $X_\lambda\to Y_\lambda$ 固有. 由命题 \bookref{B.3.4}(1), 存在指标 $\mu\in\Lambda$ 和 $X_\mu$ 上的凝聚层 $F_\mu$, 它沿 $X\to X_\mu$ 拉回为 $F$. 对 $\lambda>\mu$, 令 $F_\lambda$ 为 $F_\mu$ 沿 $X_\lambda\to X_\mu$ 的拉回. 由命题 \bookref{B.3.4}(3), 对 $\lambda\gg0$, $F_\lambda$ 在 $Y_\lambda$ 上平坦. 现在对 $\lambda\gg0$ 的数据 $X_\lambda\to Y_\lambda$ 和 $F_\lambda$ 应用 Noether 情形的上同调与基变换, 再沿 $Y\to Y_\lambda$ 基变换, 即得 $X\to Y$ 和 $F$ 的相同性质. 另见 \cite[III$_2$.$7.7.5$, III$_2$.$7.7.10$, III$_2$.$7.8.4$]{EGA}, \cite[定理 $12.11$]{Har77} 和 \cite[定理 $25.1.6$]{Vak17}.
\end{proof}
\begin{exercise}[简单而有益的练习]\label{src:A.6.10}
设 $f:X\to Y$ 为 Noether 概形之间的固有态射. 对在 $Y$ 上平坦的凝聚层 $F$, 下列条件等价:
\begin{enumerate}[(1)]
\item 对所有 $y\in Y$ 和 $i>0$, $\mathrm H^i(X_y,F_y)=0$;
\item 对所有 $i>0$, $\mathrm R^if_*F=0$, 且 $f_*F$ 为构造与 $Y$ 上基变换交换的向量丛.
\end{enumerate}
\end{exercise}
